\scanpage{8}
122\hfill J.-P. \textsc{Serre}

loi d'opération de $\operatorname{Gal}_K$ continue et semi-linéaire; on a
\[
s(cx)=s(c)s(x)\quad\text{si}\quad s\in\operatorname{Gal}_K,\ c\in C,\ x\in X.
\]
Si $i\in\mathbf Z$, notons $X^i$ le sous-ensemble de $X$ formé des éléments $x$ tels que
\[
sx=\chi(s)^i x\quad\text{pour tout}\quad s\in\operatorname{Gal}_K;
\]
c'est un $K$-espace vectoriel. Posons $X(i)=X^i\otimes_K C$. L'injection $X^i\to X$ se prolonge en une application $C$-linéaire
\[
\varepsilon_i:X(i)\to X.
\]
\textsc{Tate} (séminaire au Collège de France) a démontré le résultat suivant:

\textbf{Proposition 4.} -- \textit{Soit $\sum X(i)$ la somme directe des $X(i)$. L'homomorphisme}
\[
\varepsilon:\sum X(i)\to X,
\]
\textit{somme des $\varepsilon_i$, est injectif.}

Rappelons brièvement la démonstration. Soit $(x_{ij})$ une base de $X^i$ sur $K$. Si $\varepsilon$ n'était pas injectif, il existerait des $c_{ij}\in C$, non tous nuls, tels que
\[
\sum c_{ij}x_{ij}=0.
\]
Parmi toutes les relations de ce genre, choisissons-en une de longueur minimum, et telle que $c_{i_0j_0}=1$ pour un couple $(i_0,j_0)$ particulier. Si $s\in\operatorname{Gal}_K$, on a
\[
\sum s(c_{ij})s(x_{ij})=0,
\]
d'où
\[
\sum s(c_{ij})\chi(s)^i x_{ij}=0.
\]
Utilisant la minimalité de la relation $(c_{ij})$, on en déduit que:
\[
s(c_{ij})\chi(s)^i=\chi(s)^{i_0}c_{ij},\quad\text{pour tout}\quad s\in\operatorname{Gal}_K.
\]
Mais, d'après \textsc{Tate} [9], \S\ 3, cette dernière relation entraîne $c_{ij}=0$ pour $i\neq i_0$, et $c_{ij}\in K$ pour $i=i_0$, contrairement au fait que les $(x_{ij})$ sont linéairement indépendants sur $K$. D'où la proposition.

La proposition précédente permet \textit{d'identifier $\sum X(i)$ à un sous-espace vectoriel de $X$}. Si ce sous-espace est égal à $X$ tout entier, nous dirons que $X$ est \textit{du type Hodge-Tate}, ou admet une décomposition de type (HT).

\textbf{Remarque.} Si $X$ est de type (HT), il en est de même de son dual $X^*$, et l'on a:
\[
X^*(i)=X(-i)^*;
\]
la vérification est immédiate. De même, tout produit tensoriel de modules galoisiens de type (HT) est de type (HT).
\margin{7}
